\section{Anhang}
\begin{lstlisting}[style=bashstyle,caption={Einrichtung eines WLAN-Access-Points auf einem Raspberry Pi 4 mit Ubuntu Server},label={lst:raspi_setup}]
#!/bin/bash
# ================================================================
# Raspberry Pi 4 + Ubuntu Server
# Einrichtung eines WLAN-Access-Points mit hostapd und dnsmasq
#
# HINWEISE:
# - Dieses Skript dient als dokumentierte Installationsanleitung.
# - Einige Schritte erfordern manuelle Kontrolle der Konfigurationsdateien.
# - Vor Änderungen immer prüfen, wie die vorhandene Konfiguration aussieht.
# - Die Netplan-Datei kann anders heißen als 01-rpi.yaml.
# ================================================================

sudo apt update && sudo apt upgrade -y

sudo apt install -y hostapd dnsmasq git gpiod libgpiod-dev

sudo reboot

sudo systemctl stop hostapd
sudo systemctl stop dnsmasq

# Netplan-Datei prüfen
ls -l /etc/netplan/
sudo nano /etc/netplan/01-rpi.yaml
# Inhalt:
#
# network:
#   version: 2
#   renderer: networkd
#
#   ethernets:
#     eth0:
#       dhcp4: true

sudo chmod 600 /etc/netplan/01-rpi.yaml
sudo netplan apply

# wlan0 mit statischer IP versehen
sudo nano /etc/systemd/network/08-wlan0.network
# Inhalt:
#
# [Match]
# Name=wlan0
#
# [Network]
# Address=192.168.4.1/24

sudo systemctl restart systemd-networkd

# dnsmasq konfigurieren
sudo mv /etc/dnsmasq.conf /etc/dnsmasq.conf.orig
sudo nano /etc/dnsmasq.conf
# Inhalt:
#
# interface=wlan0
# bind-interfaces
# dhcp-range=192.168.4.20,192.168.4.100,255.255.255.0,12h
# domain-needed
# bogus-priv

# hostapd konfigurieren
sudo nano /etc/hostapd/hostapd.conf
# Inhalt:
#
# interface=wlan0
# driver=nl80211
# ssid=raspberry
# hw_mode=g
# channel=7
#
# auth_algs=1
# wpa=2
# wpa_passphrase=raspberry
# wpa_key_mgmt=WPA-PSK
# rsn_pairwise=CCMP
#
# country_code=DE
# ieee80211n=1

sudo nano /etc/default/hostapd
# Inhalt:
#
# DAEMON_CONF="/etc/hostapd/hostapd.conf"

# Dienste aktivieren
sudo systemctl unmask hostapd
sudo systemctl enable hostapd dnsmasq

sudo systemctl start hostapd
sudo systemctl start dnsmasq

# Status prüfen
systemctl status hostapd
systemctl status dnsmasq

# GPIO-Controller prüfen
gpioinfo

# Erwartete Ergebnisse:
# - hostapd: active (running)
# - dnsmasq: active (running)
# - gpioinfo liefert die erwarteten GPIO-Chips
\end{lstlisting}